\chapter{Lower bound}
\label{chap:lower_bound}

This chapter proves Theorem~3 of the paper (Section~3, proof in Appendix~C): on a hard
instance $\mathcal M_0$, every $(\eps,\delta)$-PAC algorithm uses in expectation at least
$\mathrm{lowerBound}(S,A,H,\beta,\eps,\delta,\Gmax(\mathcal M_0))$ episodes
(Theorem~\ref{thm:lower_bound}, Definition~\ref{def:lower_bound_const}). The hard instances are
those of \cite{domingues2021episodic}, with the leaf transition probabilities of the paper: a
waiting state $s_w$, an $A$-ary tree of $S - 3$ states, and two absorbing states $s_g$
(rewarding from the step $\tilde H$ on) and $s_b$. An episode of a deterministic policy
realizes exactly one \emph{triple} $u = (\text{exit step}, \text{leaf}, \text{leaf action})$, and
the instance $\mathcal M_u$ differs from $\mathcal M_0$ only in the probability of reaching $s_g$
after the triple $u$ ($p_+$ instead of $p_-$). The family thus behaves like a bandit with
$|\mathcal U| = \bar H L A = \Theta(SAH)$ arms, and the change of measure of
\cite[Lemma~1]{kaufmann2016complexity} (Theorem~\ref{thm:change_of_measure}, the paper's
Lemma~17) gives the bound.

What differs from the paper (deviations~5, 7 and~8, and the corrections recorded in
Remarks~\ref{rem:lb_condition_a}, \ref{rem:lb_params}, \ref{rem:lb_kl_paper}
and~\ref{rem:lb_constant}):
\begin{enumerate}
  \item[(i)] the tree is not required to be full: its nodes are the first $S - 3$ nodes of the
        infinite $A$-ary tree in breadth-first order, its leaves the childless nodes, and an
        action whose child is missing leads to the first child (Definition~\ref{def:hard_tree},
        Lemma~\ref{lem:tree_leaves}); the leaf action is played on arrival at the leaf, at a step
        that depends on the depth of the leaf;
  \item[(ii)] the PAC hypothesis is the one of Definition~\ref{def:entropic_pac} for the reward
        function $r_0$ of the hard instances only (deviation~7), and the bad event of the PAC
        property is the strict inequality $V^* - V^{\hat\pi} > \eps$, so the gap $\Delta = p_+ - p_-$
        is chosen so that a policy missing the triple is \emph{strictly} more than
        $\eps$-suboptimal (Definition~\ref{def:hard_params});
  \item[(iii)] Condition~A is stated with the effective horizon $H' = H - \lfloor H/3\rfloor - d$
        of the construction and with a margin ($p_+ \le (1 + p_-)/2$): the paper's condition,
        stated with $H$, does not imply admissibility with $H'$, and its bound on the binary
        divergence $\klbin(p_-, p_+)$ is false near the boundary of its condition
        (Remark~\ref{rem:lb_kl_paper});
  \item[(iv)] the change of measure is applied in LML generality to the two stationary
        environments $\mathrm{statesEnv}(\mathcal M_0)$, $\mathrm{statesEnv}(\mathcal M_u)$,
        whose actions are the policies (Lemma~\ref{lem:hard_change_of_measure});
  \item[(v)] the constant $1/1650$ becomes $1/5000$ (the proof gives $1/3456$).
\end{enumerate}
Throughout, $S := \abs{\Ss} \ge 6$ and $A := \abs{\As} \ge 2$ are the numbers of states and
actions, $\beta \ne 0$, $\eps > 0$, and $q := e^{|\beta|\eps} > 1$.

\section{The tree}
\label{sec:lb_tree}

\begin{definition}[Depth of the tree; Assumption~1]\label{def:tree_depth}
\lean{Essakine2026Tight.treeDepth}\leanok
For integers $S \ge 6$ and $A \ge 2$ let $N := (S - 3)(A - 1) + 1$ and
\[
  d := \lceil \log_A N \rceil = \big\lceil \log_A\big((S-3)(A-1) + 1\big)\big\rceil .
\]
\emph{Assumption~1} of the paper is $H \ge 3d$.
\end{definition}

\textbf{Lean remark.} \texttt{treeDepth (S A : ℕ) : ℕ := Nat.clog A ((S - 3) * (A - 1) + 1)}
(Mathlib's \texttt{Nat.clog b n} is $\lceil\log_b n\rceil$).

\begin{lemma}[Properties of the depth]\label{lem:tree_depth_props}
\uses{def:tree_depth}
For $S \ge 6$ and $A \ge 2$: $d \ge 2$, $A^{d-1} < N \le A^d$, and consequently
\[
  \frac{A^{d-1} - 1}{A - 1} < S - 3 \le \frac{A^d - 1}{A - 1} .
\]
Since $\sum_{j=0}^{k} A^j = (A^{k+1} - 1)/(A - 1)$ is the number of nodes of depth at most $k$ of
the infinite $A$-ary tree, this says that the full tree of depth $d - 2$ has fewer than $S - 3$
nodes and that the full tree of depth $d - 1$ has at least $S - 3$ nodes.
\end{lemma}

\begin{proof}
\uses{def:tree_depth}
$N \ge 3(A - 1) + 1 = 3A - 2 > A$ because $A \ge 2$, so $\log_A N > 1$ and $d \ge 2$. The
ceiling gives $d - 1 < \log_A N \le d$, i.e.\ $A^{d-1} < N \le A^d$ (Mathlib:
\texttt{Nat.pow\_pred\_clog\_lt\_self} for $N \ge 2$ and \texttt{Nat.le\_pow\_clog}). Subtracting
$1$ and dividing by $A - 1 \ge 1$, using $N - 1 = (S-3)(A-1)$, gives the displayed inequalities.
\end{proof}

\begin{definition}[The tree]\label{def:hard_tree}
\uses{def:tree_depth, lem:tree_depth_props}
The tree states are $x_0, x_1, \dots, x_{S-4}$ ($S - 3$ nodes), $x_0$ being the root, numbered in
the breadth-first order of the infinite $A$-ary tree. For $j \ge 1$ the \emph{parent} of $x_j$ is
$x_{\mathrm{par}(j)}$ with $\mathrm{par}(j) := \lfloor (j - 1)/A \rfloor$, and the \emph{depth} is
$\mathrm{dep}(x_0) := 0$, $\mathrm{dep}(x_j) := \mathrm{dep}(x_{\mathrm{par}(j)}) + 1$; equivalently
$\mathrm{dep}(x_j) = \lfloor \log_A((A-1)j + 1)\rfloor$ is the $k$ with
$(A^k - 1)/(A-1) \le j < (A^{k+1} - 1)/(A-1)$. The \emph{children} of $x_i$ are the nodes
$x_{Ai + a}$, $a \in [A] = \{1, \dots, A\}$, with $Ai + a \le S - 4$; note that
$\mathrm{par}(Ai + a) = i$ for $a \in [A]$. The node $x_i$ is \emph{internal} if it has at least
one child, i.e.\ $Ai + 1 \le S - 4$, and a \emph{leaf} otherwise; $\mathcal L$ denotes the set of
leaves and $L := \abs{\mathcal L}$. For an internal node $x_i$ and $a \in [A]$ the \emph{child
function} is
\[
  \mathrm{child}(x_i, a) := \begin{cases} x_{Ai + a} & \text{if } Ai + a \le S - 4, \\ x_{Ai + 1} & \text{otherwise (the first child),}\end{cases}
\]
so that $\mathrm{dep}(\mathrm{child}(x, a)) = \mathrm{dep}(x) + 1$ for every internal $x$ and every $a$.
\end{definition}

\textbf{Lean remark.} The state type of the hard instances has $S$ elements: an inductive type
\texttt{HardState S} with constructors \texttt{w}, \texttt{g}, \texttt{b} and
\texttt{node (j : Fin (S - 3))} (or \texttt{Fin S} with $s_w, s_g, s_b = 0, 1, 2$ and
$x_j = j + 3$), with \texttt{Fintype.card (HardState S) = S}. Actions are \texttt{Fin A}, the
paper's $a \in [A]$ being \texttt{a.val + 1}. The depth is best \emph{defined} by the closed
form \texttt{Nat.log A ((A - 1) * j + 1)} and the recursion proved as a lemma; the child
function and leaf predicate are index arithmetic (\texttt{A * i + a + 1 ≤ S - 4}).

\begin{lemma}[Leaves]\label{lem:tree_leaves}
\uses{def:hard_tree, lem:tree_depth_props}
For $S \ge 6$ and $A \ge 2$:
\begin{enumerate}
  \item[(i)] the internal nodes are $x_0, \dots, x_{I-1}$ with $I := \lceil (S-4)/A \rceil$, every
        internal node except possibly $x_{I-1}$ has all its $A$ children, and
        $L = (S - 3) - \lceil (S-4)/A \rceil$;
  \item[(ii)] $L \ge (S-3)/2 \ge S/4$ (in particular $L \ge S/6$), and $L \ge 2$;
  \item[(iii)] $\mathrm{dep}(x_j) \le d - 1$ for every $j \le S - 4$;
  \item[(iv)] for every $j \ge 1$, $a_j := j - A\,\mathrm{par}(j) \in [A]$ and
        $\mathrm{child}(x_{\mathrm{par}(j)}, a_j) = x_j$; hence for every node $x$ there is a
        sequence of $\mathrm{dep}(x)$ actions leading from $x_0$ to $x$ through the child function.
\end{enumerate}
\end{lemma}

\begin{proof}
\uses{def:hard_tree, lem:tree_depth_props}
(i) $x_i$ is internal iff $Ai + 1 \le S - 4$ iff $i \le (S-5)/A$ iff
$i \le \lfloor (S-5)/A \rfloor = \lceil (S-4)/A \rceil - 1$ (for integers $m \ge 1$,
$\lfloor (m-1)/A \rfloor + 1 = \lceil m/A \rceil$). It has all its children iff $Ai + A \le S - 4$
iff $i \le (S-4)/A - 1$, which holds for $i \le I - 2 < (S-4)/A - 1$. The leaves are the
remaining $S - 3 - I$ nodes.

(ii) $\lceil (S-4)/A \rceil \le \lceil (S-4)/2 \rceil \le (S-3)/2$ (for $S$ even
$\lceil (S-4)/2\rceil = (S-4)/2$, for $S$ odd it is $(S-3)/2$), so
$L \ge (S-3) - (S-3)/2 = (S-3)/2$; and $(S-3)/2 \ge S/4$ iff $S \ge 6$. As $L$ is an integer
and $(S-3)/2 \ge 3/2$, $L \ge 2$.

(iii) By strong induction on $j$, $\mathrm{dep}(x_j) = k$ iff
$(A^k - 1)/(A-1) \le j < (A^{k+1}-1)/(A-1)$: for $j = 0$ both sides say $k = 0$; for $j \ge 1$ with
$i = \mathrm{par}(j)$, i.e.\ $Ai + 1 \le j \le Ai + A$, if $(A^{k-1}-1)/(A-1) \le i < (A^k-1)/(A-1)$
then $j \ge A(A^{k-1}-1)/(A-1) + 1 = (A^k - 1)/(A-1)$ and
$j \le A i + A \le A\big((A^k-1)/(A-1) - 1\big) + A = (A^{k+1} - A)/(A-1) < (A^{k+1}-1)/(A-1)$.
Since $j \le S - 4 < S - 3 \le (A^d - 1)/(A-1)$ (Lemma~\ref{lem:tree_depth_props}), the depth of
$x_j$ is at most $d - 1$.

(iv) $j - 1 - A\lfloor (j-1)/A \rfloor \in \{0, \dots, A-1\}$, so $a_j \in [A]$, and
$A\,\mathrm{par}(j) + a_j = j \le S - 4$, so $x_{\mathrm{par}(j)}$ is internal and
$\mathrm{child}(x_{\mathrm{par}(j)}, a_j) = x_j$. Iterating along the ancestors
$x_j, x_{\mathrm{par}(j)}, \dots, x_0$ (a chain of length $\mathrm{dep}(x_j)$) gives the action
sequence.
\end{proof}

One checks in the same way that every leaf has depth $d - 2$ or $d - 1$; this is not needed.

\section{Parameters of the hard instances}
\label{sec:lb_params}

\begin{definition}[Condition~A]\label{def:condition_a}
\lean{Essakine2026Tight.ConditionA, Essakine2026Tight.effHorizon}\leanok
Let $\beta \ne 0$, $\eps > 0$, let $H' \ge 1$ be an integer (the \emph{effective horizon},
$H' = H - \lfloor H/3\rfloor - d$ in Definition~\ref{def:hard_params}), and put
$c := e^{|\beta| H'} - 1 > 0$ and $q := e^{|\beta|\eps} > 1$. \emph{Condition~A} for $(\beta, H', \eps)$
is
\[
  q - 1 \le \frac{c\,(2c+1)}{4\,(3c+2)} \quad (\beta > 0), \qquad\qquad
  q - 1 \le \frac{c}{10c + 8} \;\Big(\text{i.e.\ } q \le \frac{11c + 8}{10c + 8}\Big) \quad (\beta < 0).
\]
\end{definition}

\textbf{Lean remark.} \texttt{ConditionA (β : ℝ) (H' : ℕ) (ε : ℝ) : Prop}, with the two
cases by \texttt{if 0 < β then \dots\ else \dots}; Theorem~\ref{thm:lower_bound} assumes it for
\texttt{H' := H - H / 3 - treeDepth S A} (natural-number division).

\begin{definition}[Parameters of the hard instances]\label{def:hard_params}
\uses{def:tree_depth}
Let $S \ge 6$, $A \ge 2$, $H \ge 3d$ (Definition~\ref{def:tree_depth}), $\beta \ne 0$ and $\eps > 0$.
Define
\[
  \bar H := \lfloor H/3 \rfloor, \qquad \tilde H := \bar H + d + 1, \qquad H' := H - \bar H - d, \qquad
  c := e^{|\beta| H'} - 1, \qquad q := e^{|\beta|\eps}, \qquad m := \frac{3c + 2}{2(c+1)} = 1 + \frac{c}{2(c+1)},
\]
and the leaf probabilities
\begin{align*}
  \beta > 0:\qquad & p_- := \frac{1}{2(c+1)}, & \Delta &:= \frac{(q - 1)(3c+2)}{c(c+1)}, & p_+ &:= p_- + \Delta ; \\
  \beta < 0:\qquad & p_- := 1 - \frac{1}{2(c+1)}, & \Delta &:= \frac{(q-1)(3c+2)}{c(c+1)(2q-1)}, & p_+ &:= p_- + \Delta .
\end{align*}
These satisfy the identities (direct computation, using $c\Delta = 2(q-1)m$ for $\beta > 0$ and
$c\Delta = 2(q-1)m/(2q-1)$ for $\beta < 0$)
\[
  \beta > 0:\quad 1 + c\,p_- = m, \quad 1 + c\,p_+ = (2q - 1)\,m; \qquad\qquad
  \beta < 0:\quad 1 + c\,(1 - p_-) = m, \quad 1 + c\,(1 - p_+) = \frac{m}{2q-1} .
\]
Under the standing hypotheses: $d \ge 2$ (Lemma~\ref{lem:tree_depth_props}), hence $H \ge 6$,
$\bar H \ge 2$, and $\bar H \le H/3$, $d \le H/3$ give $H' \ge H/3 \ge 2$; thus $c > 0$,
$\tilde H = H - H' + 1 \le H - 1$ and $\bar H + d = H - H' \le H - 2$. Only $(\beta, \eps, H')$ enter
$c, q, m, p_\pm, \Delta$.
\end{definition}

\textbf{Lean remark.} \texttt{hardPMinus β H' : ℝ}, \texttt{hardDelta β H' ε : ℝ},
\texttt{hardPPlus β H' ε : ℝ} as functions of $(\beta, H', \eps)$ (case split on the sign of $\beta$),
and \texttt{effHorizon S A H := H - H / 3 - treeDepth S A}, \texttt{waitHorizon H := H / 3},
\texttt{rewardStart S A H := H / 3 + treeDepth S A + 1}.

\begin{remark}[Design of the parameters; comparison with the paper]\label{rem:lb_params}
By Lemma~\ref{lem:hard_exp_value} the exponential value of a policy $\pi$ in an instance is
$1 + c\,p(\pi)$ ($\beta > 0$) or $e^{\beta H'}(1 + c(1 - p(\pi)))$ ($\beta < 0$), where
$p(\pi) \in \{p_-, p_+\}$ is its probability of reaching $s_g$. The identities of
Definition~\ref{def:hard_params} say that the ratio of the exponential values of a policy with
$p_+$ and of one with $p_-$ is $2q - 1$ (resp.\ $1/(2q-1)$), so that the entropic gap between them
is $|\beta|^{-1}\log(2q - 1) > |\beta|^{-1}\log q = \eps$ (Lemma~\ref{lem:hard_eps_suboptimal}).
The paper chooses $\Delta$ so that a policy with success probability $p_- + \Delta/2$ (a
\emph{randomized} policy realizing the triple with probability $1/2$) is exactly
$\eps$-suboptimal: $\beta > 0$: $(1 + c(p_- + \Delta))/(1 + c(p_- + \Delta/2)) = q$, i.e.\
$\Delta = (q-1)(3c+2)/(c(c+1)(2 - q))$, which is our $\Delta$ times $1/(2-q)$; $\beta < 0$: the
analogous equation has the solution $(q-1)(3c+2)/(c(c+1)(2q-1))$, which is our $\Delta$, while
the paper writes $(q-1)(3c+2)/(c(c+1)(q - 1/2))$, twice the solution of its own equation.
Policies being deterministic here (Definition~\ref{def:policy}), only the gap between $p_+$ and
$p_-$ matters; it must exceed $\eps$ strictly since the bad event of
Definition~\ref{def:entropic_pac} is $V^* - V^{\hat\pi} > \eps$, and our choice is the
paper's design with the midpoint policy $p_- + \Delta/2$ made exactly $\eps$-suboptimal with
respect to the baseline $p_-$ (indeed $1 + c(p_- + \Delta/2) = q\,m$ for $\beta > 0$).
\end{remark}

\begin{lemma}[Admissibility of the parameters]\label{lem:hard_params_valid}
\uses{def:hard_params, def:condition_a}
Let $\beta \ne 0$, $\eps > 0$, $H' \ge 1$, and let $c, q, m, p_-, \Delta, p_+$ be as in
Definition~\ref{def:hard_params}. Then $c > 0$, $q > 1$, $\Delta > 0$, $0 < p_- < 1$,
$1 - p_- = (2c+1)/(2(c+1))$ if $\beta > 0$ and $1 - p_- = 1/(2(c+1))$ if $\beta < 0$, and
Condition~A for $(\beta, H', \eps)$ is equivalent to $\Delta \le (1 - p_-)/2$. Under Condition~A:
\[
  0 < p_- < p_+ \le \frac{1 + p_-}{2} < 1 \qquad\text{and}\qquad p_+(1 - p_+) \ge \frac{1}{8(c+1)} .
\]
\end{lemma}

\begin{proof}
\uses{def:hard_params, def:condition_a}
$c = e^{|\beta|H'} - 1 > 0$ since $|\beta| H' > 0$, and $q > 1$ since $|\beta|\eps > 0$; hence
$\Delta > 0$ (all factors are positive; $2q - 1 > 1$). $p_- = 1/(2(c+1)) \in (0, 1/2]$ for
$\beta > 0$ and $p_- = 1 - 1/(2(c+1)) \in [1/2, 1)$ for $\beta < 0$; the values of $1 - p_-$ follow.

\emph{Equivalence.} For $\beta > 0$, $(1-p_-)/2 = (2c+1)/(4(c+1))$ and
\[
  \Delta \le \frac{2c+1}{4(c+1)} \iff \frac{(q-1)(3c+2)}{c(c+1)} \le \frac{2c+1}{4(c+1)}
  \iff 4(q-1)(3c+2) \le c(2c+1) \iff q - 1 \le \frac{c(2c+1)}{4(3c+2)} .
\]
For $\beta < 0$, $(1-p_-)/2 = 1/(4(c+1))$ and
\[
  \Delta \le \frac{1}{4(c+1)} \iff 4(q-1)(3c+2) \le c(2q-1) \iff q(12c + 8) - 12c - 8 \le 2cq - c
  \iff q(10c + 8) \le 11c + 8 \iff q - 1 \le \frac{c}{10c+8} .
\]

\emph{Consequences.} $p_+ > p_-$ since $\Delta > 0$, and $p_+ = p_- + \Delta \le p_- + (1-p_-)/2 = (1+p_-)/2 < 1$
since $p_- < 1$. For the product: $p_+ \ge p_-$ and $1 - p_+ \ge (1 - p_-)/2$. If $\beta > 0$,
$p_- = 1/(2(c+1))$ and $(1-p_-)/2 = (2c+1)/(4(c+1)) \ge 1/4$ (as $2c + 1 \ge c + 1$), so
$p_+(1-p_+) \ge 1/(8(c+1))$. If $\beta < 0$, $p_- \ge 1/2$ and $(1-p_-)/2 = 1/(4(c+1))$, so again
$p_+(1-p_+) \ge 1/(8(c+1))$.
\end{proof}

\begin{remark}[On Condition~A]\label{rem:lb_condition_a}
\emph{The effective horizon.} The paper states Condition~A with $c = e^{|\beta|H} - 1$ while its
construction uses $c' = e^{|\beta|H'} - 1$ with $H' \le H$. The admissibility thresholds
$\varphi_+(c) := c(2c+1)/(4(3c+2))$ and $\varphi_-(c) := c/(10c+8)$ of
Definition~\ref{def:condition_a} are increasing in $c$ (the derivative of $c(2c+1)/(3c+2)$ is
$(6c^2 + 8c + 2)/(3c+2)^2 > 0$, that of $c/(10c+8)$ is $8/(10c+8)^2 > 0$), and so are the
paper's ($4(c+1)^2/(2c^2+7c+4)$ has derivative of the sign of $12c^2 + 16c + 4 > 0$; the other
one is $\varphi_-$). Since $c' \le c$, a condition stated with $H$ is \emph{weaker} than the same
condition with $H'$ and does not imply admissibility ($p_+ < 1$) for the construction: this is
the monotonicity issue of the outline, and the reason why Condition~A is stated with $H'$. As
$H' \ge H/3$ and the thresholds are increasing, Condition~A for $(\beta, \lceil H/3 \rceil, \eps)$,
or the simple sufficient condition $e^{|\beta|\eps} - 1 \le \min\{c, 1\}/20$ (both signs; for
$c \le 1$, $\varphi_+(c) \ge c/20$ and $\varphi_-(c) \ge c/18$, and for $c \ge 1$,
$\varphi_+(c) \ge 3/20$, $\varphi_-(c) \ge 1/18$), implies it.

\emph{The margin.} The paper's Condition~A is exactly the condition $p_+ < 1$ for its own
parameters (for $\beta > 0$, $p_+ < 1$ iff $q < 4(c+1)^2/(2c^2+7c+4)$, and the cap $16/13$ is used
in its bound on $\klbin(p_-,p_+)$; for $\beta < 0$, $p_+ < 1$ iff $q < (11c+8)/(10c+8)$, and its
cap $11/10$ is redundant since $(11c+8)/(10c+8) < 11/10$). Ours requires the margin
$p_+ \le (1 + p_-)/2$, which bounds $p_+(1 - p_+)$ from below (Lemma~\ref{lem:hard_params_valid})
and is what the divergence bound of Lemma~\ref{lem:hard_kl_params} needs; without a margin that
bound is false (Remark~\ref{rem:lb_kl_paper}). For $\beta < 0$ our condition coincides with the
paper's (with $H'$ for $H$); for $\beta > 0$ neither implies the other in general, ours being
more demanding for $c \le 3.1$ and less demanding for $c \ge 3.2$ (no cap at $16/13$).
\end{remark}

\section{The hard instances}
\label{sec:lb_instances}

\begin{definition}[The hard instances]\label{def:hard_mdp}
\uses{def:hard_tree, def:hard_params, def:episodic_mdp, def:reward_fn, lem:tree_leaves}
Let $S \ge 6$, $A \ge 2$, $H \ge 3d$, $\beta \ne 0$, $\eps > 0$ and the parameters of
Definition~\ref{def:hard_params}. The state space is
$\Ss := \{s_w, s_g, s_b\} \cup \{x_0, \dots, x_{S-4}\}$ ($S$ states: the waiting state, the good
and the bad absorbing states, and the tree), the action space is $\As = [A]$ with the fixed
\emph{waiting action} $a_w := 1$, the horizon is $H$ and the initial state is $s_1 := s_w$. The set
of \emph{triples} is
\[
  \mathcal U := [\bar H] \times \mathcal L \times \As, \qquad \abs{\mathcal U} = \bar H\, L\, A,
\]
and the \emph{arrival step} of $u = (h_e, \ell^*, a^*) \in \mathcal U$ is
$h_u := h_e + 1 + \mathrm{dep}(\ell^*) \le \bar H + d \le H - 2$ (Lemma~\ref{lem:tree_leaves}~(iii),
Definition~\ref{def:hard_params}). The map $u \mapsto (h_u, \ell^*, a^*)$ is injective on
$\mathcal U$ ($h_e = h_u - 1 - \mathrm{dep}(\ell^*)$). The \emph{leaf success probabilities} are,
for $h \in [H]$, $\ell \in \mathcal L$, $a \in \As$:
\[
  p^{0}_h(\ell, a) := p_- ; \qquad
  p^{u}_h(\ell, a) := \begin{cases} p_+ & \text{if } (h, \ell, a) = (h_u, \ell^*, a^*), \\ p_- & \text{otherwise.}\end{cases}
\]
For $\mathcal M \in \{\mathcal M_0\} \cup \{\mathcal M_u : u \in \mathcal U\}$, with the
corresponding $p^{\mathcal M}_h$, the transition kernel at step $h \in [H]$ is
\begin{align*}
  p_h(\cdot \mid s_w, a) &:= \delta_{s_w} \text{ if } a = a_w \text{ and } h < \bar H, \qquad \delta_{x_0} \text{ otherwise}; \\
  p_h(\cdot \mid x, a) &:= \delta_{\mathrm{child}(x, a)} \quad \text{for an internal node } x; \\
  p_h(\cdot \mid \ell, a) &:= p^{\mathcal M}_h(\ell,a)\,\delta_{s_g} + \big(1 - p^{\mathcal M}_h(\ell,a)\big)\,\delta_{s_b} \quad \text{for } \ell \in \mathcal L; \\
  p_h(\cdot \mid s, a) &:= \delta_s \quad \text{for } s \in \{s_g, s_b\} .
\end{align*}
The reward function (the same for all instances) is
\[
  r_0 : \quad r_h(s, a) := \indic\{s = s_g\}\,\indic\{h \ge \tilde H\} \in \{0, 1\},
\]
deterministic and with values in $[0,1]$, so that $\mathcal M_0, \mathcal M_u \in \mathcal M_{r_0}$
(Definition~\ref{def:reward_fn}). Thus $\mathcal M_u$ differs from $\mathcal M_0$ at exactly one
(step, state, action), namely $(h_u, \ell^*, a^*)$, where the probability of $s_g$ is $p_+$
instead of $p_-$. The definition requires $p_\pm \in [0,1]$, which holds under Condition~A
(Lemma~\ref{lem:hard_params_valid}).
\end{definition}

\textbf{Lean remark.} \texttt{hardMDP S A H β ε (u : Option (HardTriple S A H)) : EpisodicMDP (HardState S) (Fin A) H}
with $\mathcal M_0 = $ \texttt{hardMDP \dots\ none}, built from a transition function
\texttt{Fin H → HardState S → Fin A → Measure (HardState S)} by \texttt{Kernel.ofFunOfCountable}
(finite discrete types) and the deterministic reward \texttt{hardReward S A H} through
\texttt{EpisodicMDP.ofDet}. The leaf transition is the Bernoulli measure of the parameter
\emph{clamped} to $[0,1]$ (\texttt{Set.projIcc}), so that the definition makes sense for all
parameters and agrees with the above under Condition~A. The triple type is
\texttt{HardTriple S A H := Fin (H / 3) × Leaf S A × Fin A}, the paper's $h_e \in [\bar H]$ being
\texttt{h.val + 1}.

\begin{definition}[Exit step, leaf, leaf action; the triple of a policy]\label{def:hard_triple}
\uses{def:hard_mdp, def:hard_tree, def:policy, lem:tree_leaves}
Let $\pi = (\pi_h)_{h \le H}$ be a policy (Definition~\ref{def:policy}). Its \emph{exit step} is
\[
  h_e(\pi) := \min\big(\{h \in [\bar H] : \pi_h(s_w) \ne a_w\} \cup \{\bar H\}\big) \in [\bar H] .
\]
Its \emph{tree walk} is $y_0 := x_0$ and, as long as $y_k$ is internal,
$y_{k+1} := \mathrm{child}\big(y_k, \pi_{h_e(\pi) + 1 + k}(y_k)\big)$; since
$\mathrm{dep}(y_k) = k$ and all depths are at most $d - 1$ (Lemma~\ref{lem:tree_leaves}~(iii)),
the walk reaches a leaf after $K(\pi) := \min\{k : y_k \in \mathcal L\} \le d - 1$ steps. The
\emph{leaf} of $\pi$ is $\ell(\pi) := y_{K(\pi)}$, its \emph{arrival step} is
\[
  h'(\pi) := h_e(\pi) + 1 + K(\pi) = h_e(\pi) + 1 + \mathrm{dep}(\ell(\pi)) \le \bar H + d \le H - 2,
\]
its \emph{leaf action} is $a(\pi) := \pi_{h'(\pi)}(\ell(\pi))$, and its \emph{triple} is
\[
  u(\pi) := \big(h_e(\pi), \ell(\pi), a(\pi)\big) \in \mathcal U .
\]
We say that $\pi$ \emph{realizes} $u \in \mathcal U$ if $u(\pi) = u$; every policy realizes
exactly one triple. The \emph{deterministic prefix} of $\pi$ is the sequence of states
$x^\pi_h := s_w$ for $1 \le h \le h_e(\pi)$ and $x^\pi_{h_e(\pi) + 1 + k} := y_k$ for
$0 \le k \le K(\pi)$ (so $x^\pi_{h'(\pi)} = \ell(\pi)$), and
$\Phi_\pi : \{0, 1\} \to \Ss^{H+1}$ is the injective map
\[
  \Phi_\pi(b)_h := x^\pi_h \ (h \le h'(\pi)), \qquad \Phi_\pi(b)_h := s_g \text{ if } b = 1,\ s_b \text{ if } b = 0 \ (h'(\pi) < h \le H + 1).
\]
In a run $(X, Y, \mathrm{out})$ of an algorithm in the episode environment of an instance
(Definition~\ref{def:episode_env}), with $\pi^{t+1} = X_t$, the \emph{number of episodes with
triple $u$} among the first $t$ is $N_u(t) := \sum_{i < t}\indic\{u(\pi^{i+1}) = u\}$, and
$N_u(\tau)$ is its value at the stopping time on $\{\tau < \infty\}$; $E_u := \{u(\mathrm{out}) = u\}$
is the event that the output realizes $u$.
\end{definition}

\textbf{Lean remark.} \texttt{exitStep π : Fin (H / 3)} by \texttt{Nat.find} on a decidable
predicate over a finite range (or a fold over \texttt{List.range}); the tree walk by recursion on
$k \le d$ (\texttt{treeWalk π k : HardState S}), \texttt{arrivalDepth π := Nat.find} of the first
leaf (bounded by $d - 1$ by Lemma~\ref{lem:tree_leaves}), \texttt{hardTriple π : HardTriple S A H},
\texttt{Realizes π u := hardTriple π = u}; \texttt{hardTraj π : Bool → Traj (HardState S) H}
is $\Phi_\pi$ (\texttt{Traj S H := Fin (H + 1) → S}). $N_u(t)$ is a \texttt{Finset.card} of a
filtered range, a measurable function of the history.

\begin{lemma}[A realizing policy exists]\label{lem:lb_realizing_exists}
\uses{def:hard_triple, lem:tree_leaves}
For every $u = (h_e, \ell^*, a^*) \in \mathcal U$ there is a policy $\pi$ with $u(\pi) = u$.
\end{lemma}

\begin{proof}
\uses{def:hard_triple, lem:tree_leaves}
Let $a_1, \dots, a_{k}$ ($k = \mathrm{dep}(\ell^*)$) be the action sequence of
Lemma~\ref{lem:tree_leaves}~(iv) leading from $x_0$ to $\ell^*$ through the nodes
$x_0 = z_0, z_1, \dots, z_k = \ell^*$. Define $\pi_h(s_w) := a_w$ for $h < h_e$, $\pi_{h_e}(s_w) := $
any action different from $a_w$ (there is one since $A \ge 2$), $\pi_{h_e + 1 + i}(z_i) := a_{i+1}$
for $0 \le i < k$, $\pi_{h_e + 1 + k}(\ell^*) := a^*$, and $\pi$ arbitrary elsewhere. Then
$h_e(\pi) = h_e$ (if $h_e < \bar H$ because $\pi_h(s_w) = a_w$ for $h < h_e$ and
$\pi_{h_e}(s_w) \ne a_w$; if $h_e = \bar H$ by the cap), the tree walk is
$y_i = z_i$, which is internal for $i < k$ and a leaf for $i = k$, so $K(\pi) = k$,
$\ell(\pi) = \ell^*$, $h'(\pi) = h_e + 1 + k$ and $a(\pi) = \pi_{h_e+1+k}(\ell^*) = a^*$.
\end{proof}

\begin{lemma}[Trajectories of the hard instances]\label{lem:hard_trajectory}
\uses{def:hard_mdp, def:hard_triple, def:step_law, def:episode_return, lem:hard_params_valid}
Let $\mathcal M \in \{\mathcal M_0\} \cup \{\mathcal M_u : u \in \mathcal U\}$ (under
Condition~A) and let $\pi$ be a policy, $h' := h'(\pi)$ and
$p := p^{\mathcal M}_{h'}(\ell(\pi), a(\pi))$. Under $\Prob^{\pi}_{(s_w, 1)}$ of $\mathcal M$
(Definition~\ref{def:step_law}), almost surely,
\begin{enumerate}
  \item[(i)] $S_h = x^\pi_h$ for $1 \le h \le h'$ (in particular $S_{h'} = \ell(\pi)$ and $A_{h'} = a(\pi)$);
  \item[(ii)] $S_{h'+1} \in \{s_g, s_b\}$, with $\Prob(S_{h'+1} = s_g) = p$, and $S_h = S_{h'+1}$ for $h' + 1 \le h \le H + 1$.
\end{enumerate}
Consequently the state sequence $(S_1, \dots, S_{H+1})$ (Definition~\ref{def:episode_return}) has
the law $\Phi_\pi \# \mathrm{Ber}(p)$, the image of the Bernoulli law of parameter $p$ by
$\Phi_\pi$; $S_{\tilde H} \in \{s_g, s_b\}$ a.s.; and the return satisfies
$R^\pi_1 = H'\,\indic\{S_{\tilde H} = s_g\}$ a.s.
\end{lemma}

\begin{proof}
\uses{lem:step_law_markov, def:hard_mdp, def:hard_triple, def:episode_return}
(i) By induction on $h$ with the Markov property (Lemma~\ref{lem:step_law_markov}), which says
that given $S_h = s$ the action is $\pi_h(s)$ and $S_{h+1} \sim p_h(\cdot \mid s, \pi_h(s))$.
$S_1 = s_w = x^\pi_1$. For $h < h_e(\pi)$: $S_h = s_w$, $\pi_h(s_w) = a_w$ (by minimality of
$h_e(\pi)$) and $h < h_e(\pi) \le \bar H$, so $S_{h+1} = s_w = x^\pi_{h+1}$. At $h = h_e(\pi)$:
$S_h = s_w$ and either $\pi_h(s_w) \ne a_w$ or $h = \bar H$; in both cases
$S_{h+1} = x_0 = y_0 = x^\pi_{h+1}$. For $h = h_e(\pi) + 1 + k$ with $k < K(\pi)$: $S_h = y_k$ is
internal and $S_{h+1} = \mathrm{child}(y_k, \pi_h(y_k)) = y_{k+1} = x^\pi_{h+1}$ (all transitions
so far are Dirac masses). (ii) At $h = h'$: $S_{h'} = \ell(\pi)$ is a leaf, the action is
$a(\pi)$ and $S_{h'+1} \sim p\,\delta_{s_g} + (1-p)\,\delta_{s_b}$; afterwards $s_g$ and $s_b$
are absorbing, so $S_h = S_{h'+1}$ for $h > h'$ (induction with the Dirac transitions).

Since $\Phi_\pi(b)$ is exactly the sequence described by (i)--(ii) with $b = \indic\{S_{h'+1} = s_g\}$,
the state sequence is $\Phi_\pi(B)$ a.s.\ with $B \sim \mathrm{Ber}(p)$, whence the law. As
$h' + 1 \le \bar H + d + 1 = \tilde H$, $S_{\tilde H} = S_{h'+1} \in \{s_g, s_b\}$. Finally
$R^\pi_1 = \sum_{h=1}^H r_h(S_h, A_h) = \sum_{h=\tilde H}^{H}\indic\{S_h = s_g\}
= (H - \tilde H + 1)\,\indic\{S_{\tilde H} = s_g\} = H'\,\indic\{S_{\tilde H} = s_g\}$, since
$S_h = S_{\tilde H}$ for $h \ge \tilde H$ and $H - \tilde H + 1 = H'$.
\end{proof}

\textbf{Lean remark.} The useful form is the law statement: for the episode kernel
\texttt{statesKernel (hardMDP \dots\ u) w : Kernel (Policy \dots) (Traj \dots)}
(Definition~\ref{def:episode_env}),
\texttt{statesKernel (hardMDP \dots) w π = (bernoulliMeasure p).map (hardTraj π)}. It follows
from Lemma~\ref{lem:step_law_markov} by induction on the steps of the deterministic prefix
(each step being a \texttt{Kernel.deterministic}); this is the one place where the concrete
kernels of Definition~\ref{def:hard_mdp} are unfolded.

\begin{lemma}[Success probabilities]\label{lem:hard_success}
\uses{lem:hard_trajectory, def:hard_triple, def:hard_mdp}
For an instance $\mathcal M$ and a policy $\pi$ let
$p^{\mathcal M}(\pi) := \Prob^{\pi}_{(s_w,1)}(S_{\tilde H} = s_g)$ under $\mathcal M$. Then
$p^{\mathcal M}(\pi) = p^{\mathcal M}_{h'(\pi)}(\ell(\pi), a(\pi))$; hence
$p^{\mathcal M_0}(\pi) = p_-$ for every $\pi$, and for $u \in \mathcal U$,
\[
  p^{\mathcal M_u}(\pi) = p_+ \text{ if } \pi \text{ realizes } u, \qquad p^{\mathcal M_u}(\pi) = p_- \text{ otherwise.}
\]
\end{lemma}

\begin{proof}
\uses{lem:hard_trajectory, def:hard_triple, def:hard_mdp}
The first claim is Lemma~\ref{lem:hard_trajectory}~(ii) ($S_{\tilde H} = S_{h'+1}$). In
$\mathcal M_0$ all leaf probabilities are $p_-$. In $\mathcal M_u$ with $u = (h_e, \ell^*, a^*)$,
$p^{u}_{h'(\pi)}(\ell(\pi), a(\pi)) = p_+$ iff $(h'(\pi), \ell(\pi), a(\pi)) = (h_u, \ell^*, a^*)$ iff
$\ell(\pi) = \ell^*$, $a(\pi) = a^*$ and $h_e(\pi) + 1 + \mathrm{dep}(\ell^*) = h_e + 1 + \mathrm{dep}(\ell^*)$,
i.e.\ iff $u(\pi) = u$ (Definition~\ref{def:hard_triple}: $h'(\pi) = h_e(\pi) + 1 + \mathrm{dep}(\ell(\pi))$).
\end{proof}

\begin{lemma}[Exponential values]\label{lem:hard_exp_value}
\uses{def:exp_value, lem:hard_trajectory, lem:hard_success, def:hard_params}
For every instance $\mathcal M$ and policy $\pi$, with $p := p^{\mathcal M}(\pi)$,
\[
  Z^\pi_1(s_w) = 1 + p\,(e^{\beta H'} - 1) = \begin{cases} 1 + c\,p & (\beta > 0), \\ e^{\beta H'}\big(1 + c\,(1 - p)\big) & (\beta < 0). \end{cases}
\]
\end{lemma}

\begin{proof}
\uses{def:exp_value, lem:hard_trajectory, lem:hard_success}
By Lemma~\ref{lem:hard_trajectory}, $R^\pi_1 = H' B$ a.s.\ with $B \sim \mathrm{Ber}(p)$, so
$Z^\pi_1(s_w) = \E[e^{\beta R^\pi_1}] = (1 - p) + p\,e^{\beta H'}$. For $\beta > 0$ this is
$1 + c p$ with $c = e^{\beta H'} - 1$. For $\beta < 0$, $e^{\beta H'} = 1/(c+1)$ and
$(1-p) + p/(c+1) = \frac{(1-p)(c+1) + p}{c+1} = e^{\beta H'}\big(1 + c(1-p)\big)$.
\end{proof}

\begin{lemma}[Optimal exponential values]\label{lem:hard_opt_value}
\uses{def:opt_value, lem:opt_exp_value_eq, lem:hard_exp_value, lem:hard_success, lem:lb_realizing_exists, def:hard_params, lem:hard_params_valid}
Under Condition~A, for every $u \in \mathcal U$, in $\mathcal M_u$:
\[
  Z^*_1(s_w) = 1 + c\,p_+ = (2q-1)\,m \quad (\beta > 0), \qquad
  Z^*_1(s_w) = e^{\beta H'}\big(1 + c(1 - p_+)\big) = e^{\beta H'}\frac{m}{2q-1} \quad (\beta < 0),
\]
attained exactly by the policies realizing $u$; and in $\mathcal M_0$, $Z^*_1(s_w) = 1 + c\,p_- = m$
($\beta > 0$), $= e^{\beta H'}(1 + c(1-p_-)) = e^{\beta H'} m$ ($\beta < 0$).
\end{lemma}

\begin{proof}
\uses{lem:opt_exp_value_eq, lem:hard_exp_value, lem:hard_success, lem:lb_realizing_exists, lem:hard_params_valid}
By Lemma~\ref{lem:opt_exp_value_eq}, $Z^*_1(s_w) = \max_\pi Z^\pi_1(s_w)$ if $\beta > 0$ and
$= \min_\pi Z^\pi_1(s_w)$ if $\beta < 0$. By Lemma~\ref{lem:hard_exp_value}, $Z^\pi_1(s_w)$ is
$1 + c\,p^{\mathcal M}(\pi)$, increasing in $p^{\mathcal M}(\pi)$ ($c > 0$), resp.\
$e^{\beta H'}(1 + c(1 - p^{\mathcal M}(\pi)))$, decreasing in $p^{\mathcal M}(\pi)$. In
$\mathcal M_u$, $p^{\mathcal M_u}(\pi) \in \{p_-, p_+\}$ with $p_+ > p_-$
(Lemma~\ref{lem:hard_params_valid}) and $p_+$ is attained exactly by the realizing policies
(Lemma~\ref{lem:hard_success}), which exist (Lemma~\ref{lem:lb_realizing_exists}); so the
extremum is the value at $p_+$. In $\mathcal M_0$ all policies have $p_-$. The closed forms
are the identities of Definition~\ref{def:hard_params}.
\end{proof}

\begin{lemma}[Missing the triple is more than $\eps$-suboptimal]\label{lem:hard_eps_suboptimal}
\uses{def:entropic_value, def:opt_value, lem:hard_opt_value, lem:hard_exp_value, lem:hard_success, lem:hard_params_valid, def:condition_a}
Under Condition~A, for every $u \in \mathcal U$ and every policy $\pi$ that does not realize
$u$, in $\mathcal M_u$,
\[
  V^*_1(s_w) - V^\pi_1(s_w) = \frac{1}{|\beta|}\log\big(2e^{|\beta|\eps} - 1\big) > \eps .
\]
Equivalently: if $V^*_1(s_w) - V^\pi_1(s_w) \le \eps$ in $\mathcal M_u$ then $\pi$ realizes $u$
(the paper's equation (hardness)).
\end{lemma}

\begin{proof}
\uses{def:entropic_value, lem:hard_opt_value, lem:hard_exp_value, lem:hard_success, lem:hard_params_valid}
By Lemma~\ref{lem:hard_success}, $p^{\mathcal M_u}(\pi) = p_-$. If $\beta > 0$,
Lemmas~\ref{lem:hard_opt_value} and~\ref{lem:hard_exp_value} give
$V^*_1(s_w) - V^\pi_1(s_w) = \beta^{-1}\log\frac{1 + c\,p_+}{1 + c\,p_-} = \beta^{-1}\log\frac{(2q-1)m}{m} = \beta^{-1}\log(2q-1)$.
If $\beta < 0$, $V^*_1(s_w) - V^\pi_1(s_w) = \beta^{-1}\log\frac{Z^*_1(s_w)}{Z^\pi_1(s_w)}
= -|\beta|^{-1}\log\frac{e^{\beta H'} m/(2q-1)}{e^{\beta H'} m} = |\beta|^{-1}\log(2q-1)$.
(All quantities are positive: $m > 0$, $2q - 1 > 1$.) Finally $2q - 1 > q$ since $q > 1$, so
$|\beta|^{-1}\log(2q-1) > |\beta|^{-1}\log q = \eps$.
\end{proof}

\begin{lemma}[Maximal return of the hard instances]\label{lem:hard_gmax}
\uses{def:gmax, lem:hard_trajectory, lem:hard_success, lem:hard_params_valid}
Under Condition~A, $\Gmax(\mathcal M_0) = H'$ (and $\Gmax(\mathcal M_u) = H'$ for every $u$);
hence $c = e^{|\beta|\Gmax(\mathcal M_0)} - 1$ and
$c^2/(c+1) = (e^{|\beta|\Gmax(\mathcal M_0)} - 1)^2/e^{|\beta|\Gmax(\mathcal M_0)}$.
\end{lemma}

\begin{proof}
\uses{def:gmax, lem:hard_trajectory, lem:hard_success, lem:hard_params_valid}
For every policy $\pi$, $R^\pi_1 = H' B$ a.s.\ with $B \sim \mathrm{Ber}(p)$ and
$p = p^{\mathcal M}(\pi) \in \{p_-, p_+\} \subset (0, 1)$ (Lemmas~\ref{lem:hard_trajectory},
\ref{lem:hard_success}, \ref{lem:hard_params_valid}). Both values $0$ and $H'$ of $R^\pi_1$ have
positive probability, so the essential supremum of $R^\pi_1$ under $\Prob^\pi_{(s_w,1)}$ is
$H'$; the supremum over $\pi$ (Definition~\ref{def:gmax}) is $H'$.
\end{proof}

\section{Change of measure}
\label{sec:lb_change_of_measure}

For a BPI algorithm $\mathcal A$ (Definition~\ref{def:bpi_alg}) and an instance $\mathcal M$, the
law of the pair (stopped history, output) $(\mathcal H_\tau, \mathrm{out})$ of $\mathcal A$ in the
episode environment of $\mathcal M$ from $s_w$ is the same for all runs (Definition~\ref{def:bpi_alg});
we write $\Prob_0$, $\Prob_u$ for these laws under $\mathcal M_0$, $\mathcal M_u$, and $\E_0$ for
the expectation under $\Prob_0$. The number of episodes $\tau$ and the counts $N_u(\tau)$ are
functions of $\mathcal H_\tau$, so $\E_0[\tau]$ and $\E_0[N_u(\tau)]$ (in $[0, \infty]$) are defined.

\begin{lemma}[PAC algorithms find the triple]\label{lem:hard_pac_event}
\uses{def:entropic_pac, def:bpi_alg, def:hard_mdp, def:hard_triple, lem:hard_eps_suboptimal, def:reward_fn}
Assume Condition~A and let $\mathcal A$ be $(\eps,\delta)$-PAC for entropic BPI with reward
function $r_0$, parameter $\beta$ and initial state $s_w$ (Definition~\ref{def:entropic_pac}).
Then for every $u \in \mathcal U$,
\[
  \Prob_u(E_u) \ge 1 - \delta, \qquad E_u = \{u(\mathrm{out}) = u\} .
\]
Moreover, in every instance, the events $(E_u)_{u \in \mathcal U}$ are pairwise disjoint; in
particular $\sum_{u \in \mathcal U}\Prob_0(E_u) \le 1$.
\end{lemma}

\begin{proof}
\uses{def:entropic_pac, lem:hard_eps_suboptimal, def:hard_triple, def:hard_mdp}
$\mathcal M_u \in \mathcal M_{r_0}$ (Definition~\ref{def:hard_mdp}), so the PAC property gives
$\Prob_u(V^*_1(s_w) - V^{\mathrm{out}}_1(s_w) > \eps) \le \delta$ (the values being those of
$\mathcal M_u$). By Lemma~\ref{lem:hard_eps_suboptimal},
$\{u(\mathrm{out}) \ne u\} \subseteq \{V^*_1(s_w) - V^{\mathrm{out}}_1(s_w) > \eps\}$, hence
$\Prob_u(E_u^c) \le \delta$. Disjointness: $u(\mathrm{out})$ is a single element of $\mathcal U$
(Definition~\ref{def:hard_triple}), so at most one $E_u$ occurs; a sum of probabilities of
pairwise disjoint events is at most $1$.
\end{proof}

\textbf{Lean remark.} The first claim is \texttt{IdentAlg.IsPAC.measureReal\_bad\_of\_isRun}
applied to the member $\mathcal M_u$ of the class \texttt{\{M // M.HasRewardFn r₀\}}, in any run
(or to \texttt{outputMeasure} directly); the second is \texttt{measure\_biUnion\_finset} for the
pairwise disjoint preimages of the singletons $\{u\}$ by \texttt{hardTriple ∘ out}.

\begin{lemma}[Divergence of one episode]\label{lem:hard_kl_episode}
\uses{def:episode_env, lem:hard_trajectory, lem:hard_success, def:klbin, lem:klbin_bernoulli, def:hard_triple}
Assume Condition~A. For an instance $\mathcal M$ and a policy $\pi$ let
$\nu_{\mathcal M}(\pi)$ be the law of the state sequence of an episode of $\pi$ in $\mathcal M$
from $s_w$, i.e.\ the feedback kernel of the episode environment
$\mathrm{statesEnv}(\mathcal M, s_w)$ evaluated at the action $\pi$
(Definition~\ref{def:episode_env}). Then for every $u \in \mathcal U$ and every $\pi$,
\[
  \KL\big(\nu_{\mathcal M_0}(\pi) \,\big\|\, \nu_{\mathcal M_u}(\pi)\big) \le \klbin(p_-, p_+) \text{ if } \pi \text{ realizes } u,
  \qquad = 0 \text{ otherwise.}
\]
\end{lemma}

\begin{proof}
\uses{lem:hard_trajectory, lem:hard_success, lem:klbin_bernoulli, def:hard_triple}
By Lemma~\ref{lem:hard_trajectory}, $\nu_{\mathcal M}(\pi) = \Phi_\pi \# \mathrm{Ber}(p^{\mathcal M}(\pi))$
with the \emph{same} map $\Phi_\pi$ for all instances (it depends only on $\pi$ and the tree).
If $\pi$ does not realize $u$, $p^{\mathcal M_0}(\pi) = p^{\mathcal M_u}(\pi) = p_-$
(Lemma~\ref{lem:hard_success}), the two laws coincide and the divergence is $0$. If $\pi$
realizes $u$, the parameters are $p_-$ and $p_+$, and the data-processing inequality for the
measurable map $\Phi_\pi$ gives
$\KL(\Phi_\pi\#\mathrm{Ber}(p_-) \| \Phi_\pi\#\mathrm{Ber}(p_+)) \le \KL(\mathrm{Ber}(p_-)\|\mathrm{Ber}(p_+)) = \klbin(p_-, p_+)$
(Lemma~\ref{lem:klbin_bernoulli}; $p_+ \in (0,1)$ by Lemma~\ref{lem:hard_params_valid}).
\end{proof}

\textbf{Lean remark.} Data processing is Mathlib's \texttt{klDiv\_map\_le}
(\texttt{ProbabilityTheory.klDiv} of the images by a measurable map). Equality holds since
$\Phi_\pi$ is injective (apply data processing to a left inverse), but only the inequality is
used.

\begin{lemma}[Change of measure on the hard instances]\label{lem:hard_change_of_measure}
\uses{thm:change_of_measure, lem:hard_kl_episode, def:episode_env, def:bpi_alg, def:hard_triple, lem:kl_output_pair, def:klbin}
Assume Condition~A. Let $\mathcal A$ be a BPI algorithm, $u \in \mathcal U$, and assume
$\E_0[\tau] < \infty$. Then
\[
  \klbin\big(\Prob_0(E_u), \Prob_u(E_u)\big) \le \E_0\big[N_u(\tau)\big]\,\klbin(p_-, p_+) .
\]
\end{lemma}

\begin{proof}
\uses{thm:change_of_measure, lem:hard_kl_episode, def:episode_env, lem:kl_output_pair, def:hard_triple}
In LML terms (deviation~8), the episode environments $\mathrm{statesEnv}(\mathcal M_0, s_w)$ and
$\mathrm{statesEnv}(\mathcal M_u, s_w)$ (Definition~\ref{def:episode_env}) are the stationary
environments $\mathrm{stationaryEnv}(\nu_{\mathcal M_0})$ and $\mathrm{stationaryEnv}(\nu_{\mathcal M_u})$
with the finite action set $\mathcal P := \{\text{policies}\}$ and feedback kernels
$\pi \mapsto \nu_{\mathcal M_0}(\pi)$, $\pi \mapsto \nu_{\mathcal M_u}(\pi)$ (the laws of the state
sequences, Lemma~\ref{lem:hard_kl_episode}): a ``bandit'' whose arms are the policies and whose
observations are the state sequences. The sampling rule of $\mathcal A$ is an LML algorithm on
this action set, its stopping time $\tau$ is a stopping time of the history filtration, a.s.\
finite under $\Prob_0$ since $\E_0[\tau] < \infty$, and its output kernel is the same under both
environments (Lemma~\ref{lem:kl_output_pair}). For $\pi \in \mathcal P$ let
$N_\pi(\tau) := \sum_{t < \tau}\indic\{X_t = \pi\}$ be the number of pulls of the arm $\pi$
before $\tau$. Theorem~\ref{thm:change_of_measure} applied to these two environments, to the
stopping time $\tau$ and to the event $E_u = \{u(\mathrm{out}) = u\}$ of $(\mathcal H_\tau, \mathrm{out})$
gives
\[
  \klbin\big(\Prob_0(E_u), \Prob_u(E_u)\big) \le \sum_{\pi \in \mathcal P}\E_0[N_\pi(\tau)]\,\KL\big(\nu_{\mathcal M_0}(\pi)\,\big\|\,\nu_{\mathcal M_u}(\pi)\big) .
\]
By Lemma~\ref{lem:hard_kl_episode} the divergence is at most $\klbin(p_-,p_+)$ when $\pi$
realizes $u$ and $0$ otherwise, so the right-hand side is at most
$\klbin(p_-,p_+)\sum_{\pi : u(\pi) = u}\E_0[N_\pi(\tau)] = \klbin(p_-,p_+)\,\E_0[N_u(\tau)]$,
since $N_u(\tau) = \sum_{\pi : u(\pi) = u} N_\pi(\tau)$ (Definition~\ref{def:hard_triple}) and
all terms are finite ($N_\pi(\tau) \le \tau$).
\end{proof}

\textbf{Lean remark.} \texttt{statesEnv M s₁ := stationaryEnv (statesKernel M s₁)} by
definition, so Theorem~\ref{thm:change_of_measure}
(\texttt{Essakine2026Tight/LeanMachineLearning/SequentialLearning/ChangeOfMeasure.lean})
applies verbatim with the action type \texttt{Policy S A H} (a \texttt{Fintype}) and the
feedback type \texttt{Traj S H}. The pull counts $N_\pi(\tau)$ are LML's counts of the rounds
with action $\pi$ before the stopping time; the sum over $\mathcal P$ is a \texttt{Finset.sum}
over \texttt{Finset.univ}, split by \texttt{Finset.sum\_filter} according to
\texttt{hardTriple π = u}. Under $\Prob_u$ the stopping time need not be finite; whatever
convention LML uses for the output on $\{\tau = \infty\}$ is inherited by $\Prob_u(E_u)$, and
Theorem~\ref{thm:change_of_measure} is stated for it.

\begin{lemma}[Divergence of the leaf parameters]\label{lem:hard_kl_params}
\uses{def:hard_params, def:condition_a, lem:hard_params_valid, lem:klbin_upper, def:klbin}
Under Condition~A, for both signs of $\beta$,
\[
  0 < \klbin(p_-, p_+) \le 72\,\frac{c+1}{c^2}\,\big(e^{|\beta|\eps} - 1\big)^2 < \infty .
\]
\end{lemma}

\begin{proof}
\uses{lem:hard_params_valid, lem:klbin_upper, lem:klbin_bernoulli}
By Lemma~\ref{lem:hard_params_valid}, $p_+ \in (0,1)$, $p_- \in (0,1)$, $p_- \ne p_+$ and
$p_+(1-p_+) \ge 1/(8(c+1))$. Positivity and finiteness follow ($\klbin(x,y) > 0$ for
$x \ne y$, $y \in (0,1)$: Gibbs' inequality for the Bernoulli laws through
Lemma~\ref{lem:klbin_bernoulli}, Mathlib \texttt{klDiv\_eq\_zero\_iff}; and it is finite since $y \in (0,1)$). By Lemma~\ref{lem:klbin_upper}
(the bound $\klbin(x,y) \le (x-y)^2/(y(1-y))$, with $y = p_+$ the \emph{second} argument),
\[
  \klbin(p_-, p_+) \le \frac{\Delta^2}{p_+(1-p_+)} \le 8(c+1)\,\Delta^2 .
\]
Next, $(3c+2)/(c+1) \le 3$, so for $\beta > 0$, $\Delta = (q-1)(3c+2)/(c(c+1)) \le 3(q-1)/c$, and for
$\beta < 0$, $\Delta = (q-1)(3c+2)/(c(c+1)(2q-1)) \le 3(q-1)/c$ as well since $2q - 1 \ge 1$. Hence
$\klbin(p_-,p_+) \le 8(c+1)\cdot 9(q-1)^2/c^2 = 72\,(c+1)(q-1)^2/c^2$.
\end{proof}

\begin{remark}[The paper's bound on the divergence]\label{rem:lb_kl_paper}
The paper bounds $\klbin(p_-,p_+)$ by $(p_+ - p_-)^2/(p_-(1-p_-))$. The elementary bound
$\klbin(x,y) \le (x-y)^2/(y(1-y))$ (Lemma~\ref{lem:klbin_upper}, from $\log t \le t - 1$) has the
\emph{second} argument in the denominator, so the paper's expression bounds $\klbin(p_+, p_-)$,
not $\klbin(p_-, p_+)$, and no bound of the form $(x - y)^2/(x(1-x))$ can hold in general since
$\klbin(x, y) \to \infty$ as $y \to 1$ for fixed $x < 1$. This matters: the paper's Condition~A
only ensures $p_+ < 1$ without margin, and its final bound
$\klbin(p_-,p_+) \le \frac{1521}{100}\frac{c+1}{c^2}(q-1)^2$ ($\beta > 0$) is false near the
boundary. For instance with $c = 1$ and $q = 1.23$ (admissible: $q < 16/13$ and
$4(c+1)^2/(2c^2+7c+4) = 16/13$), the paper's parameters are $p_- = 1/4$,
$\Delta = 5\cdot 0.23/(2\cdot 0.77) \approx 0.7468$, $p_+ \approx 0.9968$, and
$\klbin(0.25, 0.99675) \approx 0.25\log(0.2508) + 0.75\log(230.8) \approx 3.7$, while the claimed
bound is $15.21\cdot 2\cdot 0.23^2 \approx 1.61$. Even granting the denominator, the paper drops
a factor $4$ ($p_-(1-p_-) = (2c+1)/(4(c+1)^2) \ge 1/(4(c+1))$, not $\ge 1/(c+1)$). For $\beta < 0$
the paper bounds $1/(q - 1/2)^2$ by $9e^{2|\beta|\eps}$ where $4$ suffices ($q \ge 1$); this
is the only source of the factor $e^{2\min\{\beta,0\}\eps}$ in its Theorem~3, which is therefore
unnecessary (it is $\le 1$ and is kept in Definition~\ref{def:lower_bound_const} only to match
the paper's statement). Our Condition~A includes the margin $p_+ \le (1+p_-)/2$, which makes
the bound $\klbin(p_-,p_+) \le \Delta^2/(p_+(1-p_+)) \le 8(c+1)\Delta^2$ valid.
\end{remark}

\begin{lemma}[The counts sum to the stopping time]\label{lem:hard_count_sum}
\uses{def:hard_triple, def:bpi_alg}
In every run of a BPI algorithm in the episode environment of an instance, on $\{\tau < \infty\}$,
$\sum_{u \in \mathcal U} N_u(\tau) = \tau$. Consequently, if $\tau < \infty$ $\Prob_0$-a.s.,
$\E_0[\tau] = \sum_{u \in \mathcal U}\E_0[N_u(\tau)]$ (in $[0,\infty]$).
\end{lemma}

\begin{proof}
\uses{def:hard_triple}
Each of the $\tau$ episodes $t < \tau$ has exactly one triple $u(\pi^{t+1})$
(Definition~\ref{def:hard_triple}), so $\sum_u N_u(\tau) = \sum_{t<\tau}\sum_u\indic\{u(\pi^{t+1}) = u\} = \tau$.
Take expectations (a finite sum of nonnegative terms).
\end{proof}

\begin{lemma}[Lower bound on the sum of divergences]\label{lem:hard_count_lower}
\uses{def:klbin, lem:klbin_lower}
Let $\mathcal U$ be a finite set with $\abs{\mathcal U} \ge 4$, $\delta \in (0, 1/16]$,
$(x_u)_{u \in \mathcal U}$ in $[0,1]$ with $\sum_u x_u \le 1$ and $(y_u)_{u \in \mathcal U}$ in
$[1 - \delta, 1]$. Then, in $[0, \infty]$,
\[
  \sum_{u \in \mathcal U}\klbin(x_u, y_u) \;\ge\; \sum_{u \in \mathcal U}\Big[(1 - x_u)\log\frac1\delta - \log 2\Big] \;\ge\; \frac{\abs{\mathcal U}}{2}\log\frac1\delta .
\]
\end{lemma}

\begin{proof}
\uses{lem:klbin_lower}
\emph{Termwise.} If $y_u < 1$, Lemma~\ref{lem:klbin_lower} gives
$\klbin(x_u,y_u) \ge (1-x_u)\log\frac{1}{1-y_u} - \log 2 \ge (1-x_u)\log\frac1\delta - \log 2$
since $1 - y_u \le \delta$ and $1 - x_u \ge 0$. If $y_u = 1$ and $x_u < 1$, $\klbin(x_u, 1) = \infty$
(Definition~\ref{def:klbin}). If $x_u = y_u = 1$, $\klbin(1,1) = 0 \ge -\log 2$.

\emph{Summing.} $\sum_u[(1-x_u)\log\frac1\delta - \log 2] = (\abs{\mathcal U} - \sum_u x_u)\log\frac1\delta - \abs{\mathcal U}\log 2
\ge (\abs{\mathcal U} - 1)\log\frac1\delta - \abs{\mathcal U}\log 2$. This is at least
$\frac{\abs{\mathcal U}}{2}\log\frac1\delta$ iff $(\frac{\abs{\mathcal U}}{2} - 1)\log\frac1\delta \ge \abs{\mathcal U}\log 2$,
which holds because $\frac{\abs{\mathcal U}}{2} - 1 \ge \frac{\abs{\mathcal U}}{4}$ (as
$\abs{\mathcal U} \ge 4$) and $\frac14\log\frac1\delta \ge \frac14\log 16 = \log 2$ (as $\delta \le 1/16$).
\end{proof}

\textbf{Lean remark.} $\klbin$ is the $\texttt{ℝ≥0∞}$-valued \texttt{klBin} of
Definition~\ref{def:klbin} (as \texttt{klBer} in the sister project, infinite when the second
argument is $1$ and the first is not); the sum is a \texttt{Finset.sum} in \texttt{ℝ≥0∞} and the
right-hand side is \texttt{ENNReal.ofReal} of a real number. The outline's version
($\abs{\mathcal U} \ge 2$, constant $1/4$) is weaker; the paper's constant $1/2$ is correct
since $\abs{\mathcal U} \ge 4$ in the application.

\section{The lower bound}
\label{sec:lb_theorem}

\begin{definition}[The lower bound]\label{def:lower_bound_const}
\lean{Essakine2026Tight.lowerBound}\leanok
\uses{def:gmax}
For integers $S, A, H \ge 1$, $\beta \ne 0$, $\eps > 0$, $\delta \in (0,1)$ and $G \ge 0$,
\[
  \mathrm{lowerBound}(S, A, H, \beta, \eps, \delta, G) := \frac{1}{5000}\,\frac{(e^{|\beta| G} - 1)^2}{e^{|\beta| G}}\,
  \frac{e^{2\min\{\beta, 0\}\eps}\, S A H}{(e^{|\beta|\eps} - 1)^2}\,\log\frac1\delta .
\]
\end{definition}

\textbf{Lean remark.} \texttt{lowerBound (S A H : ℕ) (β ε δ G : ℝ) : ℝ} in
\texttt{Essakine2026Tight/EV2026/Constants.lean}.

\begin{theorem}[Theorem~3, lower bound]\label{thm:lower_bound}
\lean{Essakine2026Tight.exists_hardMDP_lowerBound_le_lintegral_stoppingTime}\leanok
\uses{def:tree_depth, def:condition_a, def:hard_params, def:hard_mdp, def:entropic_pac, def:bpi_alg, def:episode_env, def:gmax, def:lower_bound_const, def:reward_fn}
Let $S \ge 6$, $A \ge 2$, $H \ge 3d$ with $d$ as in Definition~\ref{def:tree_depth}, $\beta \ne 0$,
$\delta \in (0, 1/16]$ and $\eps > 0$ satisfying Condition~A for $(\beta, H', \eps)$ with
$H' := H - \lfloor H/3 \rfloor - d$ (Definition~\ref{def:condition_a}). There are a reward
function $r_0$ with values in $[0,1]$, a finite-horizon MDP $\mathcal M_0$ with $S$ states, $A$
actions, horizon $H$ and reward function $r_0$, and an initial state $s_1$, such that every BPI
algorithm $\mathcal A$ which is $(\eps,\delta)$-PAC for entropic BPI with reward function $r_0$,
parameter $\beta$ and initial state $s_1$ (Definition~\ref{def:entropic_pac}) satisfies, in every
run in the episode environment of $\mathcal M_0$ from $s_1$,
\[
  \E_0[\tau] \;\ge\; \mathrm{lowerBound}\big(S, A, H, \beta, \eps, \delta, \Gmax(\mathcal M_0)\big) .
\]
More precisely the proof gives $\E_0[\tau] \ge \frac{1}{3456}\frac{(e^{|\beta|\Gmax(\mathcal M_0)} - 1)^2}{e^{|\beta|\Gmax(\mathcal M_0)}}\frac{SAH}{(e^{|\beta|\eps}-1)^2}\log\frac1\delta$.
\end{theorem}

\begin{proof}
\uses{lem:hard_params_valid, lem:hard_gmax, lem:hard_pac_event, lem:hard_change_of_measure, lem:hard_kl_params, lem:hard_count_sum, lem:hard_count_lower, lem:tree_leaves, lem:tree_depth_props, def:hard_triple}
Take $r_0$, $\mathcal M_0$, the family $(\mathcal M_u)_{u \in \mathcal U}$ and $s_1 = s_w$ of
Definition~\ref{def:hard_mdp}, which is admissible by Lemma~\ref{lem:hard_params_valid}. Let
$\mathcal A$ be $(\eps,\delta)$-PAC for $r_0$, fix a run in the episode environment of
$\mathcal M_0$, and write $B := \mathrm{lowerBound}(S,A,H,\beta,\eps,\delta,\Gmax(\mathcal M_0)) < \infty$.

\emph{Step 0.} If $\E_0[\tau] = \infty$ there is nothing to prove. Assume $\E_0[\tau] < \infty$;
then $\tau < \infty$ $\Prob_0$-a.s.

\emph{Step 1: change of measure.} By Lemma~\ref{lem:hard_change_of_measure}, for every
$u \in \mathcal U$, $\klbin(\Prob_0(E_u), \Prob_u(E_u)) \le \E_0[N_u(\tau)]\,\klbin(p_-,p_+)$. Summing
over $u$ and using Lemma~\ref{lem:hard_count_sum},
\[
  \sum_{u \in \mathcal U}\klbin\big(\Prob_0(E_u), \Prob_u(E_u)\big) \le \klbin(p_-,p_+)\sum_{u}\E_0[N_u(\tau)] = \klbin(p_-,p_+)\,\E_0[\tau] .
\]

\emph{Step 2: the PAC property.} By Lemma~\ref{lem:hard_pac_event}, $x_u := \Prob_0(E_u)$
satisfies $\sum_u x_u \le 1$ and $y_u := \Prob_u(E_u) \ge 1 - \delta$. Moreover
$\abs{\mathcal U} = \bar H L A \ge 1\cdot 2\cdot 2 = 4$ ($\bar H \ge 1$, $L \ge 2$ by
Lemma~\ref{lem:tree_leaves}, $A \ge 2$), so Lemma~\ref{lem:hard_count_lower} gives
$\sum_u \klbin(x_u, y_u) \ge \frac{\abs{\mathcal U}}{2}\log\frac1\delta$. With Step~1 and
$0 < \klbin(p_-,p_+) < \infty$ (Lemma~\ref{lem:hard_kl_params}),
\[
  \E_0[\tau] \ge \frac{\abs{\mathcal U}\log(1/\delta)}{2\,\klbin(p_-,p_+)} \ge \frac{\abs{\mathcal U}\log(1/\delta)}{2\cdot 72}\,\frac{c^2}{(c+1)(q-1)^2}
  = \frac{\abs{\mathcal U}}{144}\,\frac{c^2}{c+1}\,\frac{\log(1/\delta)}{(e^{|\beta|\eps}-1)^2} .
\]

\emph{Step 3: the number of triples.} $\bar H = \lfloor H/3 \rfloor \ge H/6$ for $H \ge 3$ (for
$H \ge 4$, $\lfloor H/3\rfloor \ge (H-2)/3 \ge H/6$; for $H = 3$, $1 \ge 1/2$), $L \ge S/4$
(Lemma~\ref{lem:tree_leaves}~(ii)), so $\abs{\mathcal U} = \bar H L A \ge SAH/24$.

\emph{Step 4: the constant.} By Lemma~\ref{lem:hard_gmax},
$c^2/(c+1) = (e^{|\beta|\Gmax(\mathcal M_0)}-1)^2/e^{|\beta|\Gmax(\mathcal M_0)}$. Hence
\[
  \E_0[\tau] \ge \frac{1}{24\cdot 144}\,\frac{(e^{|\beta|\Gmax(\mathcal M_0)}-1)^2}{e^{|\beta|\Gmax(\mathcal M_0)}}\,\frac{SAH}{(e^{|\beta|\eps}-1)^2}\log\frac1\delta ,
\]
$24 \cdot 144 = 3456 \le 5000$, and $e^{2\min\{\beta,0\}\eps} \le 1$; this is $\E_0[\tau] \ge B$.
\end{proof}

\textbf{Lean remark.} \texttt{Essakine2026Tight.exists\_hardMDP\_lowerBound\_le\_lintegral\_stoppingTime}:
for the state type \texttt{HardState S} (cardinality $S$) and action type \texttt{Fin A}, the
statement is
$\exists\, r_0\ M_0\ s_1,\ M_0.\texttt{HasRewardFn}\ r_0 \wedge (\forall a\, h\, s, r_0\, h\, s\, a \in [0,1]) \wedge
\forall\, \mathcal A,\ \texttt{IsEntropicPAC}\ \mathcal A\ r_0\ \beta\ \eps\ \delta\ s_1 \to \forall \text{ run},\
\texttt{ENNReal.ofReal}\ (\mathrm{lowerBound}\ \dots\ (\Gmax\ M_0\ s_1)) \le \int \tau\, d\Prob$,
with $\tau : \Omega \to \texttt{ℕ∞}$ cast to \texttt{ℝ≥0∞} (so $\E_0[\tau] = \infty$ is allowed and
Step~0 is the case \texttt{lintegral = ⊤}). The hard MDP $M_0$ is \texttt{hardMDP S A H β ε none}
and the family $\mathcal M_u$ only appears in the proof. Condition~A is a hypothesis on
$(\beta, \texttt{effHorizon S A H}, \eps)$.

\begin{remark}[On the constant and the hypotheses]\label{rem:lb_constant}
The proof gives the constant $1/3456 = 1/(2\cdot 72\cdot 24)$: $2$ from
Lemma~\ref{lem:hard_count_lower}, $72$ from Lemma~\ref{lem:hard_kl_params} and $24 = 6\cdot 4$
from $\bar H \ge H/6$, $L \ge S/4$. Under Assumption~1 one even has $H \ge 6$ and
$\lfloor H/3\rfloor \ge H/4$, giving $1/2304$; we state $1/5000$ (deviation~5) to leave slack.
The paper's $1/1650$ rests on its erroneous divergence bound (Remark~\ref{rem:lb_kl_paper}),
on $L \ge S/4$ for a full tree (which our Lemma~\ref{lem:tree_leaves} recovers for the
breadth-first tree) and on $\bar H = H/3$ taken as an integer. The hypotheses $\delta \le 1/16$
(used only in Lemma~\ref{lem:hard_count_lower}) and $H \ge 3d$ (used to fit the waiting phase,
the tree and a rewarding phase of length $H' \ge H/3$ into the horizon) are the paper's; the
paper's Theorem~3 writes $\delta \in (0, 1/16)$ in the main text and $\delta \le 1/16$ in the
appendix, and the proof only needs $\log(1/\delta) \ge 4\log 2$. The lower bound is stated for
algorithms that are PAC for the single reward function $r_0$ (deviation~7), a weaker hypothesis
than PAC for every reward function, hence a stronger conclusion.
\end{remark}
